\section{Stopping a search without losing the possibilities it has not examined}
\label{sim:temporal-frontier}

A correction may have reached the manager before a payment decision, or afterward. An interval for the delivery time can preserve both possibilities without assigning either a probability. If the source record contains several corrections, public releases and private deliveries, the exact temporal calculation may exhaust its work allowance before checking every combination. The unfinished calculation still has a useful result, provided it keeps the completed conclusions together with an exact description of what remains unexamined.

The conservative earlier method returned the whole original clock domain whenever its search stopped. That answer preserved uncertainty but discarded useful completed proof work. A resumable calculation can retain more. It must save the partition it has already established, including every unfinished part. Saving only the branches that produced interesting financial views would change the question on the next run.

\subsection{Each split must cover the entire timing domain}

Consider one cell $D$ of admissible event times under the supplied chronology and access assumptions. For the next source revision, let $g_1,\ldots,g_m$ be its eligibility conditions. A condition may require that publication or actor access occurred no later than the decision, or that ingestion occurred no later than the system's reconstruction date. The revision is eligible only if all these conditions hold.

A complete split keeps one eligible cell and one cell for each possible first failed condition:
\begin{align}
 D_{\mathrm{yes}} &= D\cap\bigcap_{j=1}^{m}g_j,\\
 D_{\mathrm{no},k} &= D\cap\bigcap_{j<k}g_j\cap g_k^{\mathrm c},
 \qquad k=1,\ldots,m.
\end{align}
Here a condition also denotes the set of time assignments satisfying it, and $g_k^{\mathrm c}$ denotes its exact complement. Every assignment in $D$ either satisfies every condition or has a unique first failure. These alternatives therefore cover $D$ and are pairwise disjoint. Their union is the original cell, not a sample of it. Some children may be empty, and a contradiction proof can establish that fact later.

The word \emph{exact} matters at a decision boundary. If $g_k$ says that access occurred at or before time $d$, its complement says strictly after $d$. Replacing that complement by \emph{at or after} would put the boundary in two children. Replacing the original condition by a strict inequality would remove the boundary from the eligible child. The selected reference retains rational event times and strict inequalities, so the preceding submicrosecond and fixed-offset examples keep their original meanings.

The search resolves one pending cell at a time. It first checks whether the cell is empty. A nonempty cell with no remaining revision to classify yields a selected or unresolved source view. Otherwise it installs the complete family of children above. Only after all those children exist may the calculation pause. This logical rule prevents a work limit from turning a complete split into a partial list of favorable branches.

\subsection{Pending does not mean feasible, and known does not mean complete}

A saved node has one of four meanings: its domain remains unchecked; a contradiction proves it empty; a witness establishes a terminal source view; or a witness supports an expanded cell whose children form the complete partition. An unchecked child need not contain a possible time assignment. It may already be inconsistent with an earlier chronology restriction, but the calculation has not yet established that. Retaining it is conservative; deleting it without checking is unsupported.

In the selected example, resolving the root installs eight children. Seven are later proved empty. With a one-node work allowance, all eight remain explicitly pending. The result therefore neither asserts eight feasible worlds nor pretends that the single eventual view is already established. The eight cells are a representation of unresolved work under known restrictions.

Once some terminal views are known, they can be displayed together with their source supports. Agreement among them still does not establish that the whole family has that view. An unexamined cell could contain a conflicting correction, a retraction or a different access history. Completion requires that no pending cell remains, including cells that eventually prove empty.

The same distinction controls policy interpretation. A policy that fails in an already established view fails for any larger family containing that view. A policy that succeeds in every established view may still fail in an unexamined one. Thus an incomplete temporal search can sometimes establish a counterexample before it can establish universal feasibility. Connecting that asymmetry to the numerical policy calculation requires an explicit treatment of the pending remainder; a saved search alone does not supply it.

\subsection{Resume the original question and preserve its proofs}

The checkpoint binds the complete source request, actor, trust declarations, timing assumptions, method identities and execution allowance. On reading it, a separate checker reconstructs each branch's constraints from the original domain and its path. It does not accept a saved mutable constraint list as a substitute. It verifies all completed witnesses or contradictions, the exact complement structure, source outcomes, support paths and work counts.

The reference uses a fixed left-first depth-first order. This makes the resumed tree directly comparable with the existing uninterrupted calculation. It also allows the checker to reject a purported completed branch placed after an earlier unfinished one. The order is a comparison convention, not a claim that this schedule is the fastest way to find useful financial conclusions.

A step allowance specifies how much new work one call may perform. A separately bound total execution cap determines when further work must stop. Reaching that cap retains the pending tree and reports a capped calculation; it does not turn unresolved cells into completed results. The test with a cap of two preserves the same original question and performs no repeated solver work when resumed at that cap. Changing the cap or the source assumptions requires an explicitly different, checked request to the continuation interface.

Across the 47 selected temporal requests, step allowances of one, two, seven and 4,096 produce exactly the same final proof trees, outcomes and source supports as uninterrupted search. Earlier certificates remain unchanged at every pause. A separate reread checks 1,060 saved checkpoints and rejects altered partition cells, proofs, source premises and completion declarations. These comparisons establish the selected finite continuation property.\footnote{The implementation, independent structural checker, 188 complete comparisons, 43 main adversarial controls and five additional saved-record controls are retained in \path{data/accounting-master/temporal-frontier-20261027/}. The main comparison makes 10,146 assertions; the saved-record audit makes 3,191. The source locks and manifests preserve the exact inputs, methods and comparison scope.}

This representation avoids a separate explicit node for every real-valued event-time assignment. It does not remove exponential growth in the number of distinguishable revision combinations. Nor does installing all children before a logical pause establish crash-safe file storage, authenticated evidence or actual actor knowledge. Larger revision families, better schedules, justified approximate reduction and durable storage require their own arguments and comparisons. The useful advance here is narrower: unfinished computation retains both what has been proved and every domain that has yet to be checked.
